%==============================================================================
% Section 6 -- Differentiable Rasterisation Pipeline and Tiling
%==============================================================================
\section{Differentiable Rasterisation Pipeline and Tiling}
\label{sec:rasterization}

\subsection{Design Goals}

The pipeline must be: (1) \textbf{real-time} ($<$10 ms/frame at 1080p),
(2) \textbf{differentiable} (gradients back to all Gaussian parameters), and
(3) \textbf{depth-correct} (front-to-back compositing). These are achieved
via a \textbf{tile-based, sorting-based CUDA rasteriser}.

%------------------------------------------------------------------------------
\subsection{Pipeline Overview}

\begin{algorithm}[H]
  \caption{3DGS Differentiable Rasterisation Pipeline}
  \label{alg:pipeline}
  \begin{algorithmic}[1]
    \Require{Gaussian set $\mathcal{S}$, camera parameters, resolution $W \times H$}
    \Ensure{Rendered image $\hat{I}$}
    \State \textbf{Preprocessing} (parallel over all $N$ Gaussians):
      project means, compute 2D covariances, evaluate SH colours, compute
      bounding boxes and depths.
    \State \textbf{Tile assignment}: divide screen into $16\times16$ tiles;
      for each Gaussian, append a 64-bit key
      $[\tau\texttt{\_id}\,\|\,z_i]$ for every overlapping tile.
    \State \textbf{GPU radix sort} by key $\Rightarrow$ depth order per tile
      ($O(M)$, $\sim10^9$ keys/s).
    \For{each tile $\tau$ \textbf{in parallel} (one CUDA block)}
      \For{each pixel $\mathbf{u}$ in $\tau$}
        \State Accumulate colour via alpha blending
               (Algorithm~\ref{alg:alpha_blend}).
      \EndFor
    \EndFor
    \State \Return{$\hat{I}$}
  \end{algorithmic}
\end{algorithm}

%------------------------------------------------------------------------------
\subsection{Tile Decomposition}

The screen is split into $16\times16$ pixel tiles (256 threads/block, a
multiple of the CUDA warp size). Pixels in the same tile share most Gaussians,
enabling cooperative shared-memory loading.

%------------------------------------------------------------------------------
\subsection{CUDA Rendering and Early Termination}

Each CUDA block loads batches of 128 Gaussians into shared memory and each
thread accumulates its pixel colour independently. A thread maintains
accumulated transmittance $T$; when $T < \epsilon = 10^{-4}$ the pixel is
\emph{saturated} and the thread exits, providing significant speed-ups for
occluded pixels.

%------------------------------------------------------------------------------
\subsection{Backward Pass}

A custom CUDA kernel re-traces the forward pass in \emph{reverse tile order}
(back-to-front), accumulating gradients
$\partial\mathcal{L}/\partial\mean'_i$,
$\partial\mathcal{L}/\partial\cov'_i$,
$\partial\mathcal{L}/\partial o_i$,
$\partial\mathcal{L}/\partial\mathbf{c}_i$
via atomic operations. World-space gradients are then obtained by autodiff
through the preprocessing pass. The overall pipeline scales as $O(N)$ in
Gaussians and $O(WH)$ in pixels.
